\section{Cutting into pants and varying twists}
\label{sec:twists}

\subsection{What data a cut curve retains}

Fix a pants decomposition with cuff lengths $a_1,\ldots,a_N$, where
$N=3g-3+n$. Cutting a simple multicurve gives disjoint essential arcs in
each pair of pants. Cutting a general immersed curve, even a filling one,
need not give disjoint arcs: self-intersections may remain inside a pair
of pants. Nor are the actual segments generally orthogonal to the cuffs.

Boundary intersection counts alone lose the relative homotopy classes of
the arcs, their cyclic joining pattern, and winding information. For simple
multicurves this information is organized by Dehn--Thurston coordinates.
For a general immersed curve one must retain a fuller itinerary. Once this
topological data is fixed, varying the metric twist is a well-defined and
tractable problem.

There are two complementary exact descriptions. One uses arc lengths as
generating functions in endpoint positions. The other uses products of
holonomy matrices. Both keep the information that intersection counts alone
discard.

\subsection{An exact local distance function}

Take two ultraparallel geodesic lifts of cuffs in $\HH^2$, at common
perpendicular distance $d$. Let $u,v$ be signed distances on the two lines
from the perpendicular feet, with compatible orientations. Then
\begin{equation}
 \boxed{
 D_d(u,v)=\arcosh\bigl(\cosh d\cosh u\cosh v-\sinh u\sinh v\bigr).}
 \label{eq:localdistance}
\end{equation}
For example, in the hyperboloid model use
\[
 p(u)=(\cosh u,\sinh u,0),\qquad
 q(v)=(\cosh d\cosh v,\sinh v,\sinh d\cosh v).
\]
Then $\cosh D=-\langle p,q\rangle$. Reversing one endpoint coordinate
changes the sign of the mixed term.

For the ordinary seam between cuffs $a,b$ of pants with third cuff $c$,
\begin{equation}
 \cosh d=\frac{\cosh(c/2)+\cosh(a/2)\cosh(b/2)}
                    {\sinh(a/2)\sinh(b/2)}.
 \label{eq:seam}
\end{equation}
For a returning or winding arc, choose the boundary lifts appropriate to
its relative homotopy class and use their common perpendicular instead.
There is no single seam length that describes all relative arc classes.

Given a fixed itinerary, cuff gluing imposes affine relations on the
endpoint variables involving the twists $\tau_i$ and integer winding
data. The length of its closed geodesic representative is obtained by
minimizing the sum of the corresponding functions~\eqref{eq:localdistance}
over the crossing positions. The first-variation equations impose smooth
tangent matching at every crossing. This provides an exact formulation
even when the resulting minimization is not elementary to evaluate.

\subsection{Holonomy and first variation}

In a basis adapted to a cuff axis, a twist displacement is represented by
\[
 T_i(\tau_i)=\diag(e^{\tau_i/2},e^{-\tau_i/2}).
\]
For a fixed curve word $\gamma$, insert these matrices into the fixed pants
holonomies. The resulting product $W_\gamma$ satisfies
\begin{equation}
 2\cosh(\ell_\gamma/2)
  =\abs{\tr W_\gamma(\mathbf a,T_1(\tau_1),\ldots,T_N(\tau_N))}.
 \label{eq:traceword}
\end{equation}
With the cuff lengths fixed, its trace is a finite Laurent polynomial in
$e^{\tau_i/2}$. Its coefficients need not all be positive in an arbitrary
choice of lifts and coordinates.

The twist derivative has the cosine formula
\begin{equation}
 \frac{\partial\ell_\gamma}{\partial\tau_i}
   =\sum_{p\in\gamma\cap\alpha_i}\cos\vartheta_p,
 \qquad
 \abs{\frac{\partial\ell_\gamma}{\partial\tau_i}}
      \le i(\gamma,\alpha_i).
 \label{eq:cosine}
\end{equation}
Angles and twist orientation are chosen consistently. Length is strictly
convex along an individual nontrivial twist when the curve crosses the
twisting cuff; the argument also applies to nonsimple geodesics
\cite{Kerckhoff}. We do not need a stronger claim about arbitrary lines
with several twist parameters varying in different signs.

These ideas overlap with several parts of Mirzakhani's work, but serve
different roles there. Dehn--Thurston estimates compare combinatorial and
geometric length for simple curves~\cite[Section 3]{MirzSimple}. The work
on general curve orbits studies asymptotically piecewise-linear length,
including filling curves and length sublevel volumes
\cite[Sections 4 and 8]{MirzGeneral}. In the basic volume recursion, the
cutting curves are simple and the kernel depends on their lengths, so the
twist integral itself is just the period $a_i$.

\subsection{Unfolding a twist lattice}

Let $D_i$ be the Dehn twist about the $i$th cuff. With a consistent sign
convention,
\[
 \ell_{D_1^{k_1}\cdots D_N^{k_N}\gamma}(\mathbf a,\boldsymbol\tau)
 =\ell_\gamma(\mathbf a,\boldsymbol\tau+\mathbf k\mathbf a).
\]
For a nonnegative function $F$, Tonelli's theorem gives
\begin{equation}
 \boxed{
 \int_{\prod_i[0,a_i)}\sum_{\mathbf k\in\ZZ^N}
 F\bigl(\ell_{D_1^{k_1}\cdots D_N^{k_N}\gamma}
                  (\mathbf a,\boldsymbol\tau)\bigr)\,\dd\boldsymbol\tau
 =\int_{\RR^N}F\bigl(\ell_\gamma(\mathbf a,\mathbf u)\bigr)\,\dd\mathbf u.}
 \label{eq:twistunfold}
\end{equation}
For signed functions, impose absolute integrability. If the curve has a
nontrivial stabilizer in this twist lattice, remove the redundant directions
or account for their multiplicity; in particular, twisting about a disjoint
cuff should not produce infinitely many copies of the same term.

A fixed curve's length is not periodic in a twist coordinate. The sum over
its twist orbit is periodic, which is exactly what makes this unfolding
work. The accompanying WP measure is $\prod_i\dd a_i\,\dd\tau_i$.
However, the positive cuff orthant times one twist box is the quotient by
this particular twist subgroup, not a fundamental domain for the full
mapping class group. A further full-moduli integration needs the appropriate
orbit unfolding.

For example, with weight $(2\cosh(\ell_\gamma/2))^{-s}$ and
$x_i=e^{u_i/2}$, the right side becomes the Mellin-type integral
\begin{equation}
 2^N\int_{(0,\infty)^N}
 \abs{\tr W_\gamma(\mathbf a,\mathbf x)}^{-s}
             \prod_i\frac{\dd x_i}{x_i},
 \label{eq:mellin}
\end{equation}
whenever it converges. This is an explicit analytic object attached to the
curve itinerary and the fixed pants lengths.

